\documentclass[Total.tex]{subfiles}

\begin{document}
	\chapter{Topological Spaces}
		\section{Definitions}
		
			\subsection{Topological Space}
			\begin{Def}
				内容...
			\end{Def}
			\subsection{Generating of Topological Spaces}
			\subsection{Bases, Subbases}
			
			\subsection{Continuous Mapping}
		\section{Constructions}
			\subsection{Quotient Space}
				\subsubsection{Quotient Topology}
				\begin{Def}
					Suppose $X$ be a set with equivalence relation $(X,\sim)$:
					\begin{itemize}
						\item Given $x\in X$, we denote $[x]$ the unique equivalence class that contains $x$;
						\item We refer to the map
						\begin{gather*}
							p:X\rightarrow X/\sim\qquad x\mapsto [x]
						\end{gather*}
						Natural projection;
						\item Let $f:X\rightarrow Y$ be a map between sets such that $f(x)=f(y)$ whenever $x\sim y$. By the quotient factorization the map
						\begin{gather*}
							g:X/\sim\rightarrow Y\qquad [x]\mapsto f(x)
						\end{gather*}
						is well-defined and it is the unique map $g:X/\sim\rightarrow Y$ such that $f=g\circ p$.
					\end{itemize}
				\end{Def}
				
				\begin{proposition}
					(Topological-Quotient Proposition)
				\end{proposition}
				\begin{proof}
					内容...
				\end{proof}
				\begin{Def}
					内容...
				\end{Def}
				\subsubsection{Open Map, Closed Map}
			\subsection{Product, Coproduct}
			\subsection{Limits and Colimits}
			\subsection{Pullback,Pushout}
				
				\subsubsection{Pullback}
				\begin{proposition}
					Let $p:X\rightarrow B$ be a continuous map of topological spaces. Let $C$ be a topological space and $f:C\rightarrow B$ be a continuous map. We consider the topological space:
					\begin{gather*}
						f^*X:=\{(c,x)\in C\times X|f(c)=p(x)\}
					\end{gather*}
					and the maps
					\begin{gather*}
						\varphi:f^*X\rightarrow C\qquad (c,x)\mapsto c\\
						\psi:f^*X\rightarrow X\qquad (c,x)\mapsto x
					\end{gather*}
					The following statements hold:
					\begin{itemize}
						\item \begin{enumerate}
							\item The maps $\varphi:f^*X\rightarrow C$ and $\psi:f^*X\rightarrow X$ are continuous;
							\item Suppose we are given a topological space $W$ and two continuous maps $\alpha:W\rightarrow C$ and $\beta:W\rightarrow X$ such that
							\begin{gather*}
								f\circ \alpha=p\circ \beta
							\end{gather*}
							The map
							\begin{gather*}
								\Phi:W\rightarrow f^*X\\
								w\mapsto (\alpha(w),\beta(w))
							\end{gather*}
							is continuous. In particular we obtain the following commutative diagram
							\begin{center}
								\begin{tikzcd}
									W\arrow[rd,"\exists!\Phi",dashrightarrow]\arrow[rrd,"\beta",bend left=10]\arrow[ddr,"\alpha",bend right=10]\\
									&f^*X\arrow[r,"\psi"]\arrow[d,"\varphi"]&X\arrow[d,"p"]\\
									&C\arrow[r,"f"]&B
								\end{tikzcd}
							\end{center}
						\end{enumerate}
						A map $\Phi:W\rightarrow f^*X$ is continuous iff the maps $\varphi\circ\Phi:W\rightarrow C$ and $\psi\circ\Phi:W\rightarrow X$ are continuous.
						\item The topological space $f^*X$ with natural maps $\varphi:f^*X\rightarrow C$ and $\psi:f^*X\rightarrow X$ is the pullback, in the category $\mathcal Top$.
					\end{itemize}
				\end{proposition}
				\begin{proof}
					\begin{itemize}
						\item 
						\begin{itemize}
							\item Consider $f^*X$ as a subspace of $W=X\times C$. Then, $\varphi$ and $\psi$ can be seen as restriction of $\pi_C,\pi_X$ on $f^*X$.
							\item 
							\item 
						\end{itemize}
						\item 
					\end{itemize}
				\end{proof}
				\begin{proposition}
					(Pullback-Properties Proposition) Let $p:X\rightarrow B$ be a continuous map of topological spaces, let $C$ be a topological space and let $f:C\rightarrow B$ be a continuous map:
					\begin{itemize}
						\item If $X$ and $C$ are compact and if $X,B,C$ are Hausdorff, then $f^*X$ is compact;
						\item Connectedness: \begin{itemize}
							\item Even if $X,B,C$ are connected, it does not necessarily follow that the pullback $f^*X$ is connected;
							\item Suppose that $X,B,C$ are connected. We also assume that for each $b\in B$ the preimage $p^{-1}(b)\subseteq X$ is connected. It does not follow $f^*X$ is connected;
						\end{itemize}
						\item If $X,C$ are Hausdorff, the pullback $f^*X$ is also Hausdorff.
					\end{itemize}
				\end{proposition}
				\begin{proof}
					内容...
				\end{proof}
				\subsubsection{Pushout}
				\begin{proposition}
					(Topological-Pushout Proposition) Let $f:A\rightarrow Y$ and $g:A\rightarrow Z$ be two continuous maps between topological spaces. We consider the topological space
					\begin{gather*}
						Y\cup_A Z:=(Y\sqcup Z)/\sim\qquad f(a)\simeq g(a),a\in A
					\end{gather*}
					and the map
					\begin{gather*}
						i:Y\rightarrow Y\sqcup_A Z\qquad j:Z\rightarrow Y\sqcup_A Z
					\end{gather*}
					The following statements hold:
					\begin{itemize}
						\item $U\subseteq Y\cup_A Z$ is open$\Leftrightarrow i^{-1}(U)\subseteq Y,j^{-1}(U)\subseteq Z$ are open.
						\item 
						\item 
					\end{itemize}
				\end{proposition}
				\begin{proof}
					内容...
				\end{proof}
				
				\begin{Def}
					内容...
				\end{Def}
				\begin{proposition}
					(Pushout-Properties Proposition) Let $f:A\rightarrow Y,g:A\rightarrow Z$ be two continuous maps betwween topolgoical spaces. We consider the corresponding pushout 
				\end{proposition}
			\subsection{Fibre Product}
			\subsection{Specialization}
				\begin{Def}
					In topology, the \textbf{specialization preoreder} is a preorder on the point sets of topological space.
					
					For \textbf{T0 Space}, this is a partial order. called \textbf{specialization order}: For topological space $(X,\tau)$,
					
					\begin{gather*}
						x\leq y\Leftrightarrow x\in \overline{\{y\}}
					\end{gather*}
				\end{Def}
				
				\begin{proposition}
					We have:
					\begin{itemize}
						\item If the space is $T_0$, the preorder is a partial order.
						\item The symmetry of specialization preorder: $x\leq y\wedge y\leq x\rightarrow x=y\Leftrightarrow R_0$ Axiom;
						\item For $T_1$ Space, we have
						\begin{gather*}
							x\leq y\Leftrightarrow x=y
						\end{gather*} 
						\item For sober space with specialization order $(X,\leq)$:
							\begin{itemize}
								\item 
								\item 
							\end{itemize}
					\end{itemize}
				\end{proposition}
				\begin{proof}
					内容...
				\end{proof}
				
				\subsubsection{Morphsims between Specialization Preorders}
					
		\section{Locally Closed Space}
			\begin{proposition}
				The following properties are equivalent:
				\begin{itemize}
					\item $E$ is the intersection of an open set and a closed set in $X$; 
					\item For each $x\in E$, there is a neighborhood $U$ of $x$ such that $E\cap U$ of $x$ such that $E\cap U$ is closed in $U$;
					\item $E$ is open in its closure $\bar{E}$;
					\item The set $\bar{E}-E$ is closed in $X$;
					\item $E$ is the difference of two closed sets in $X$;
					\item $E$ is difference of two open sets in $X$;
				\end{itemize}
			\end{proposition}
			\begin{proof}
				\begin{itemize}
					\item (1)$\Rightarrow$ (2)
					\item (2)$\Rightarrow$ (3)
					\item (3)$\Rightarrow$ (4)
					\item (4)$\Rightarrow$ (5)
					\item (5)$\Rightarrow$ (6)
					\item (6)$\Rightarrow$ (1)
				\end{itemize}
			\end{proof}
		\section{Noetherian Space}
		
			\begin{Def}
				\begin{itemize}
					\item A topological space is called \textbf{Noetherian} if the descending chain condition holds for closed subsets of $X$. 
					
					\item A topological space is called \textbf{locally Noetherian} if every point has a neighborhood which is Noetherian.
				\end{itemize}
			\end{Def}
			
			\begin{lemma}
				Let $X$ be a topological space:
				
				\begin{itemize}
					\item Any subset of $X$ with the induced topology is Noetherian;
					\item The space $X$ has finitely many irreducible components;
					\item Each irreducible component of $X$ contains a nonempty open of $X$;
				\end{itemize}
			\end{lemma}
			
			\begin{proof}
				\begin{itemize}
					\item 
					\item 
					\item 
				\end{itemize}
			\end{proof}
			
			\begin{lemma}
				(Noetherian Property via Continuous Mapping) Let $f:X\rightarrow Y$ be a continuous map of topological spaces:
				
				\begin{itemize}
					\item If $X$ is Noetherian, then $f(X)$ is Noetherian;
					\item If $X$ is locally Noetherian and $f$ open, then $f(X)$ is locally Noetherian;
				\end{itemize}
			\end{lemma}
			
			\begin{proof}
				内容...
			\end{proof}
			
			\begin{lemma}
				Let $X$ be a topological space and $X_i\subseteq X;i=1,\dots,n$. If each $X_i$ is Noetherian, then $\cup X_i$ is Noetherian(with induced topology).
			\end{lemma}
			
			\begin{proof}
				
			\end{proof}
			
			\begin{lemma}
				Let $X$ be a locally Noetherian topological space, then $X$ is locally connected.
			\end{lemma}
			
			\begin{proof}
				内容...
			\end{proof}
		\section{Elementary Knowledge in Dimension}
		
			\subsection{Krull Dimension}
				
				\begin{Def}
					Let $X$ be a topological space:
					
					\begin{itemize}
						\item A chain of \textbf{irreducible closed subsets} of $X$ is a sequence
						
						\begin{gather*}
							Z_0\subset\dots \subset Z_n\subset X
						\end{gather*} 
						
						with $Z_i$ 
						
						\item 
						\item 
						\item 
					\end{itemize}
				\end{Def}
				
				
			
			\subsection{Codimension and Catenary Maps}
			
				\begin{Def}
					Let $X$ be a topological space. Let $Y\subseteq X$ be an \textbf{irreducible} closed subset. The \textbf{codimension} of $Y$ in $X$ is the supremum of lenghts $e$ of chains
					
					\begin{gather*}
						Y=Y_0\subset \dots\subset Y_e\subset X
					\end{gather*}
				\end{Def}
				
				\begin{lemma}
					Let $X$ be a topological space, and $Y\subseteq X$ is an irreducible closed subset. And $U\subseteq X$ is an open subset such that $Y\cap U$ is nonempty, then
					
					\begin{gather*}
						codim(Y,X)=codim(Y\cap U,U)
					\end{gather*}
				\end{lemma}
				
				\begin{proof}
					内容...
				\end{proof}
			\subsection{Catenary Space}
			
				\begin{Def}
					内容...
				\end{Def}
		\section{Jacobson Spaces}
			
			\begin{Def}
				Let $X$ be a topological space, and $X_0$ be the set of \textbf{closed points} in $X$(i.e. $\overline{\{x\}}=\{x\}$). We say that: $X$ is \textbf{Jacobson} if every closed subset $Z\subseteq X$ is the closure of $Z\cap X_0$. 
				
			
			\end{Def}
			
				\begin{lemma}
				A topological space $X$ is Jacobson iff every locally closed subset of $X$ has a point closed in $X$.
			\end{lemma}
			
			\begin{proof}
				内容...
			\end{proof}
			
			\textbf{Remark}
			
			\begin{lemma}
				Let $X$ be a topological space and $X_0$ be a set of closed points in $X$. Suppose that for every $x\in U\cap Z$, we have that
				
				\begin{gather*}
					内容...
				\end{gather*}
			\end{lemma}
			
			\begin{proof}
				内容...
			\end{proof}
			
		\section{Limit of Spces}
	\chapter{Filter,Nets}
		\section{Convergence of Filters}
			\subsection{Filter}
				\begin{Def}
					Let $E$ be a set: A family of subsets of $E$ satisfies
					\begin{itemize}
						\item $\varnothing\notin\mathcal{F}$;(F1)
						\item $A,B\in \mathcal{F}\Rightarrow A\cap B\in \mathcal{F}$;(F2)
						\item $B\in \mathcal{F},B\subseteq A\Rightarrow A\in \mathcal{F}$;(F3)
					\end{itemize}
				\end{Def}
				
				\begin{example}
					Consider a sequence of points $x_1,\dots,x_n\in E$:
					\begin{Def}
						The \textbf{associated filter} is the family of all subsets of $E$ such that: The subset of $E$ contains all $x_i,i\in \NN$ except possibly a finite number of items.
					\end{Def}
				\end{example}
				
				\subsubsection{Basis of Filter}
				
				\begin{Def}
					A family $\mathfrak{B}$ is a \textbf{basis of filter} $\mathcal{F}$ if:
					\begin{itemize}
						\item $\mathfrak{B}\subseteq \mathcal{F}$;(BF1)
						\item Every subset of $E$ belonging to $\mathcal{F}$ contains some subsets of $E$ belongs to $\mathfrak{B}$:
						\begin{gather*}
							S\in \mathcal{F}\Rightarrow \exists T\in \mathfrak{B}, T\subseteq S
						\end{gather*}(BF2)
					\end{itemize}
				\end{Def}
				
				\begin{example}
					Filter of neighborhoods of $0$ on real line:
					\begin{gather*}
						\mathcal{F}=\{(-a,a)|a>0\}
					\end{gather*}
				\end{example}
				
				\begin{example}
					Consider a sequence with $S=\{x_i|i\in \NN\}$, then the filter given by
					\begin{gather*}
						S_n:=\{x_n,x_{n+1},\dots\},S_0\supseteq S_1\supseteq \dots	 
					\end{gather*}
					is a basis of $\mathcal{F}$.
				\end{example}
				
				\textbf{Question} Given some family of subsets $\mathcal{A}$. Is there a filter $\mathcal{F}$ having basis $\mathcal{A}$ as a basis?
				
				\begin{lemma}
					Given two filter $\mathcal{F}_1,\mathcal{F}_2$ with basis of $\mathfrak{B}$. Then, $\mathcal{F}_1=\mathcal{F}_2$. It is uniquely determined.
				\end{lemma}
				\begin{proof}
					Suppose $S\in \mathcal{F}_1$, then there exists some $T\in \mathfrak{B}$ such that $T\subseteq S$. And then, $S\in \mathcal{F}_2$.
				\end{proof}
				
				So finally, this leads the question to: Is $\mathcal{F}$ a filter? (F2) equivalent to:
				\begin{itemize}
					\item The intersection of any two sets beloing to $\mathcal{A}$, contains a set which belongs to $\mathcal{A}$.(BF)
				\end{itemize}
				This is weaker than (F2). Thus we may state:
				\begin{center}
					A basis of filter on $E$ is a family of \underline{nonempty subsets} of $E$ satisfying condition (BF).
				\end{center}
				The filter is uniquely determined by (BF2)
				
				\subsubsection{Comparison of Filters}
				Consider the order relation induced by $(P(E),\subseteq)$:
				\begin{itemize}
					\item Finer;
					\item less fine;
				\end{itemize}
				
				\subsubsection{Topology}
				The definition of topology can be rewritten as following:
				\begin{itemize}
					\item $S\in \mathcal{F}(x)\Rightarrow x\in S$;(N1)
					\item $U\in \mathcal{F}(x)\Rightarrow \exists V\in \mathcal{F}(x),\forall y\in V,S\in \mathcal{F}(y)$.(N2)
				\end{itemize}
				Such $\mathcal{F}(x)$ is called the \textbf{filter of neighborhoods} of the point $x$.
			\subsection{Convergence with Filter}
			\begin{Def}
				\begin{itemize}
					\item Define: Filter $\mathfrak{F}$ converges to the point $x\Leftrightarrow$ Every neighborhood of $x$ belongs to $\mathfrak{F}$.
					
					In other word: $\mathcal{F}$ is finer filter than neighborhoods of $x$.
					\item Sequence $(x_n)$ converges to $x$ if its associated filter converges to $x$.
				\end{itemize}
			\end{Def}
			\subsection{Continuity}
				Then: The continuity can be rewritten as the following form
				\begin{gather*}
					内容...
				\end{gather*}
		\section{Moore-Smith Convergence}
			\subsection{Directed Set, Nets}
				\begin{Def}
					内容...
				\end{Def}
	\chapter{Connectedness and Irreducibility}
	
		\section{Connectedness}
		
		\section{Irreduciblity}
	
	\chapter{Compactness}
		\section{Compactness}
		
			\subsection{Compact Spaces}
			
			\begin{Def}
				\begin{itemize}
					\item We say: A space is \textbf{(quasi-)comapct} if every open covering $X=\cup U_i$ has a finite subcovering;
					\item A continuous map $f:X\rightarrow Y$ is \textbf{(quasi-)comapct} iff the inverse image $f^{-1}(V)$ is compact for any compact $V\subseteq Y$;
					\item A subset $Z\subseteq X$ is called \textbf{retrocompact} if for every compact open $U\subseteq X$, $Z\cap U$ is compact;
				\end{itemize}
			\end{Def}
			
			\subsection{Compact Maps}
			
		\section{Locally Compact Spaces}
		\section{Constructible Sets}
		
			\begin{Def}
				Let $X$ be a topological space and $E\subseteq X$ is a subset:
				
				\begin{itemize}
					\item We say $E$ is \textbf{constructible} in $E$ iff $E$ is a finite union of subsets of the form
					
					\begin{gather*}
						U\cap V^C
					\end{gather*}
					
					where $U,V$ are both open and retrocompact;
					
					\item We say $E$ is \textbf{locally constructible} if for there exists an open covering $X=\cup V_i$ such that each $E\cap V_i$ is constructible in $V_i$;
				\end{itemize}
			\end{Def}
			
			\begin{lemma}
				The collection of constructible sets is closed under:
				
				\begin{itemize}
					\item Finite Union;
					\item Finite Intersection;
					\item Compelement;
				\end{itemize}
			\end{lemma}
			
			\begin{proof}
				\begin{itemize}
					\item 
					\item 
					\item 
				\end{itemize}
			\end{proof}
		\section{Partition of Unity}
		
	\chapter{Separation Axioms}
	
	\chapter{Conuntable Axioms}
	
	\chapter{Function Spaces and Metrization Problem}
	
	\chapter{Baire Spaces}
	
	\chapter{Dimension}
		
\end{document}